\documentclass[10pt]{article}
\usepackage[margin=20mm]{geometry}
\usepackage{amsmath,amssymb,mathtools,fontspec,unicode-math,xurl,hyperref}
\IfFontExistsTF{DejaVu Serif}{\setmainfont{DejaVu Serif}}{\setmainfont{Latin Modern Roman}}
\setmathfont{latinmodern-math.otf}
\setlength{\parindent}{0pt}\setlength{\parskip}{6pt}\sloppy
\title{Scientific initial review: M25-0004
Fundamental groups and the Milnor conjecture}
\author{}
\date{}
\begin{document}
\maketitle
\section*{Version, scope, and assessment}
This review concerns only the manifest-bound manuscript of E. Bruè, A. Naber and D. Semola, arXiv:2303.15347v3 (20 November 2023), 56 PDF pages, with its bound 2497-line `source.tex`. The input PDF SHA256 is\par
`3dfc6be8569357835d843ae24072756ae969547c4a418def304c952ed64c44e6`.\par
Page numbers below refer to this manuscript.\par
The checked arXiv record still lists v3 as its latest version. The publisher records a revision on 25 April 2024 and publication in Annals of Mathematics 201 (2025), 225–289, DOI 10.4007/annals.2025.201.1.4. Thus the journal revision postdates this input. No journal-final PDF or equivalence proof was acquired. A bounded primary-source search did not identify an article-specific erratum; this is not evidence that none exists. A separate six-dimensional sequel is a distinct paper, not a replacement version or correction receipt for this one.\par
All 56 pages were read. The principal dependency is the equivariant positive-Ricci path (Section 6), its constant-volume refinement and neck (Section 7), and the action-extension gluing (Section 8), feeding the direct limit in Section 3. This review identifies three Major derivation defects and three Minor scope/description defects in v3. Major here means a demonstrated invalid step on a core dependency; it does not mean that the main existence theorem has been disproved. The independent calculations, static-source locations, coverage, version records, build evidence and page checks are supplied separately.\par
\section*{F01: the fixed connection is not Yang–Mills throughout the proposed deformation}
Major. The locations are Lemma 6.3(ii), p. 29; (67)–(75), pp. 29–31; and the harmonic-extension assertion in Section 6.6, p. 34. The corresponding source ranges are 1072–1178 and 1257–1265. The assertion that the displayed family has zero mixed Ricci curvature is false. It is not a consequence of the two diagonal Christoffel identities in (70).\par
\subsection*{Direct computation from the actual ansatz}
Suppress the parameter \(t\), and use the orthonormal frame in (67): \[\begin{gathered} T=f^{-1}(U_1+kV_1),\quad X_2=U_2,\quad X_3=U_3,\\
 W_1=a^{-1}(-kU_1+V_1),\quad
 W_2=(ab)^{-1}V_2,\quad W_3=(ab)^{-1}V_3.\end{gathered}\] Here \(a,b,f>0\). Using the paper’s simultaneous bracket-sign convention, the bracket \([T,W_2]\) is \(2kW_3/f\), not \(2kX_3/f\) as printed in (68). A simultaneous sign reversal for right-invariant frames does not change any Ricci calculation below.\par
The exact mixed component obtained from the Koszul formula is \[\operatorname{Ric}(T,W_1)=\frac{2kf}{(1+k^2)^2}
       \left(\frac{1}{a^3b^4}-a\right).
 \tag{R1}\] For completeness, in a unimodular invariant orthonormal frame one can compute this component as the sum of the bracket-norm, Killing-form and output-bracket terms. The first two cross terms cancel here. The pair \([X_2,X_3]\), including its reverse order, contributes \(-2kfa/(1+k^2)^2\); the pair \([W_2,W_3]\) contributes \(2kf/((1+k^2)^2a^3b^4)\). Their sum is (R1). The independently executed script instead forms all brackets, all Christoffel coefficients, the curvature tensor and its contraction, thereby checking the same identity without copying (74) or (75).\par
Take \(k=1,\ a=1,\ b=1/2,\ f=1/10\). This is a smooth positive invariant metric of exactly the displayed form, with constant-length circle fibers and the fixed connection \(\eta_0\). It gives \[\operatorname{Ric}(T,W_1)=\frac34,\qquad
 \operatorname{Ric}(T,T)=\frac{17}{200},\qquad
 \operatorname{Ric}(W_1,W_1)=\frac{17}{2}.\] Since the circle fibers have constant length, the mixed Ricci formula in Proposition 4.4 makes nonzero mixed curvature equivalent, up to the stated nonzero scalar and sign convention, to nonzero divergence of the curvature form. Thus (74), and the zero assertion in (75), fail.\par
\subsection*{Why the endpoint requirements do not remove the problem}
For \(k\ne0\), (R1) vanishes precisely when \(ab=1\). Rounding the quotient \(S^3\) fibers requires \(b=1/(1+k^2)\), by the three lengths displayed on p. 30. Consequently zero mixed curvature at that round endpoint would force \(a=1+k^2\). Direct computation of the base Ricci component gives \[\operatorname{Ric}_{N}(X_2,X_2)=4-\frac{2k^2a^2}{(1+k^2)^2}.
 \tag{R2}\] At this endpoint it equals \(4-2k^2\), which is negative for \(|k|\ge2\). For example \(k=2,a=5,b=1/5\) gives \(-4\). Therefore the proposed fixed-connection ansatz cannot simultaneously provide the stated Yang–Mills property, round fibers and positive base Ricci curvature for these values of \(k\). This conclusion concerns the constructed ansatz; it is not a classification of all possible metrics in the lemma.\par
There is also an immediate incompatible condition at source lines 1102–1104: \(b\) is required to be nondecreasing from \(1/\sqrt{1+k^2}\) to the smaller number \(1/(1+k^2)\). Reversing this inequality repairs that particular sign, but does not repair (R1) or the endpoint obstruction. Nor does correcting the isolated bracket index repair the Yang–Mills assertion.\par
\subsection*{Dependency, counterevidence, and required repair}
Section 6.6 chooses the unique harmonic curvature representative and asserts that it agrees with the already constructed curvature for all \(t\in[0,1/2]\). This agreement uses precisely the false divergence claim. One therefore needs a different connection/path construction, or a proved transition to harmonic representatives with full mixed-curvature control. The dependency is Lemma 6.3 \(\to\) Section 6.6 \(\to\) Theorem 6.1 \(\to\) Lemma 7.2 and Proposition 7.1 \(\to\) the induction for Theorem 1.1. The issue is consequently Major.\par
The calculation also supplies substantial counterevidence against an overstated conclusion: the \(T,W_1\) block has determinant \[\operatorname{Ric}(T,T)\operatorname{Ric}(W_1,W_1)-\operatorname{Ric}(T,W_1)^2
   =\frac{4f^2}{(1+k^2)^2a^2b^4}>0.\] In the explicit \(k=1\) example the other four diagonal entries are \(699/200,699/200,198/25,198/25\), and all other cross entries vanish. The total metric in that example has positive Ricci curvature despite not being Yang–Mills. This finding invalidates the printed proof step, not the possibility of a corrected positive-Ricci deformation.\par
\section*{F02: the neck Ricci formula omits the second variation of its cross section}
Major. The location is (128), p. 45, source lines 1694–1699, and its use in (135), p. 46. For \(G=dr^2+h(r)^2g(r)\) with constant cross-section volume form, the printed tangential formula has no \(g''\) term. Such a term cannot be eliminated by constant volume.\par
An explicit test uses \(S^3\times S^3\) with \[h(r)=1,\qquad
 g(r)=e^{2c(r-1)^2}g_{S^3}\oplus e^{-2c(r-1)^2}g_{S^3},\qquad c>0.
 \tag{R3}\] This is smooth on an open interval about \(1\), has constant volume form, and is invariant under every Hopf action under discussion. At \(r=1\), \(g'=0\), while \(g''\) has blocks \(4c\,g_{S^3}\) and \(-4c\,g_{S^3}\). Formula (128) predicts Ricci curvature \(2\) in a unit direction of the first sphere. The ordinary doubly warped formula, or a direct shape-operator calculation, gives instead \[\operatorname{Ric}_G(X_1,X_1)=2-2c.
 \tag{R4}\] For \(c=2\) these values are \(2\) and \(-2\), respectively. All prerequisites for the general formula are satisfied; the test does not purport to be the authors’ particular final neck.\par
Here is the missing mechanism explicitly. Put \(q=h^2g\) and \(K=q'/2\). For tangential indices, \[\operatorname{Ric}^G_{ij}=\operatorname{Ric}^{q}_{ij}-K'_{ij}
       -(\operatorname{tr}_q K)K_{ij}+2K_{i\ell}q^{\ell m}K_{mj}.\] With \(\operatorname{tr}_g g'=0\) and cross-section dimension six this becomes \[\begin{aligned}
 \operatorname{Ric}^G_{ij}={}&\operatorname{Ric}^g_{ij}
 -(hh''+5(h')^2)g_{ij}-3hh'g'_{ij}\notag\\
 &-\frac{h^2}{2}g''_{ij}
   +\frac{h^2}{2}g'_{i\ell}g^{\ell m}g'_{mj}.
 \end{aligned}\tag{R5}\] The specialization \(h=1,g'=0\) already proves the discrepancy, without depending on the other coefficients in (R5).\par
The estimate (135) does not bound this missing term. A repair must bound the second derivative of the compact smooth metric family and use \[g''=\ddot{\widetilde g}\,(t')^2+\dot{\widetilde g}\,t''.\] Slow reparametrization offers a plausible way to make these errors small; compact smoothness supplies finite higher derivative bounds. These are counterevidence to claiming the neck cannot be constructed. They do not make the printed estimate valid, and this review has not certified a replacement set of simultaneous constants. The defect is on the Proposition 7.1 neck construction used by Theorem 1.1, so it is Major.\par
\section*{F03: the volume normalization does not justify the Moser equation}
Major. Equations (118)–(120), p. 43, source lines 1618–1650, assert constant volume after defining \[\widetilde g_r=a_r\varphi_r^*\widehat g_r,\qquad
 a_r=\widetilde a\,\operatorname{Vol}(\widehat g_r)^{-1/6}.\] With the displayed tensor multiplication in dimension six, \[\operatorname{Vol}(a_r\widehat g_r)
  =a_r^3\operatorname{Vol}(\widehat g_r)
  =\widetilde a^{\,3}\operatorname{Vol}(\widehat g_r)^{1/2},\] which is not generally constant. A diffeomorphism cannot alter the total volume. Moreover the subsequent definition is \(\nu_r=dv_{\widehat g_r}\), so even replacing the exponent alone would leave the stated density equation unnormalized.\par
A complete countertest to this normalization step is \(\widehat g_r=c(r)g_0\), with \(c>0\) smooth, equal to \(1\) near both endpoints and nonconstant in the interior. For \(k=0\) this is already an admissible positive-Ricci equivariant path with the specified standard endpoints; one may also insert a scalar bump in any admissible path. Its volume is \(V_0c(r)^3\). The proposed scaled tensor has volume \(\widetilde a^3V_0^{1/2}c(r)^{3/2}\), which varies.\par
The obstruction to (120) is exact, not a numerical concern: the weighted integral of its left side is zero by integration by parts, whereas the required integral of its right side is \(-\frac{d}{dr}\int\nu_r\). This is nonzero where the above volume changes. Thus the stated solvability condition need not hold.\par
A concrete algebraic repair is to use the tensor factor \(a_r=\widetilde a\,\operatorname{Vol}(\widehat g_r)^{-1/3}\) and then define \(\nu_r=dv_{a_r\widehat g_r}\). Equivalently one can consistently use a length scale and its square in every metric formula. With this repair the integral obstruction vanishes, and the usual invariant weighted elliptic equation yields an equivariant Moser isotopy. This does not address F01 or F02. The incorrect normalization is itself used to obtain Lemma 7.2(3), an explicit prerequisite for the neck calculation, and is Major under the required core-derivation standard despite its elementary repair.\par
\section*{Three Minor defects with explicit limiting counterevidence}
\subsection*{F04: coprime weights do not ensure freeness at the zero section}
The location is p. 11, Section 2.2.1, and Propositions 3.1(4) and 5.1(4), pp. 16 and 24; source lines 423, 569 and 840. In the constructed product model, take \((a,b)=(2,1)\). These weights are coprime and \(a\ne0\). At any \((u,0)\in S^3\times\mathbb R^4\), the nonidentity element \(\theta=\pi\) acts by \((e^{2i\pi}u,e^{i\pi}0)=(u,0)\). More generally the zero-section stabilizer has order \(|a|\). Away from the zero section, coprimality does eliminate the simultaneous stabilizer. Global freeness for this model requires \(|a|=1\).\par
The actual inductive construction uses weights \((1,k)\), which pass this test, and its cyclic subactions are free. Restricting the overgeneral assertion to \(|a|=1\) leaves that mechanism available. This is therefore Minor, not a counterexample to the free action needed by the theorem.\par
\subsection*{F05: the universal metric-cone description contradicts the stated cylinder limit}
The final paragraph of Section 2.2, p. 14 (source line 494), says that every tangent at infinity is a metric cone. Section 2.2.7, p. 13, equation (22), and the introduction explicitly include a limit \(\mathbb R^3\times S^1\) when \(k_j\to\infty\). A full metric cone contracts to its vertex by radial contraction and is contractible. The cylinder has fundamental group \(\mathbb Z\), so it is not isometric or homeomorphic to any such cone, regardless of basepoint. Both assertions cannot be retained simultaneously.\par
A bounded sequence of indices can avoid this particular limit; the statement in question was not restricted to that case, and Example 2.3 allows unbounded indices. Delete or restrict the universal description, and distinguish a general pointed blow-down limit from a metric cone. This is a descriptive contradiction in the outline, not an obstruction to the inductive manifold construction, hence Minor.\par
\subsection*{F06: the mapping-class argument silently omits orientation reversal}
Equation (173) and the four-generator conclusion on p. 54 are written for all of \(\operatorname{Diff}(S^3\times S^3)\) with image in \(\mathrm{SL}(2,\mathbb Z)\). The diffeomorphism \(\sigma(x,y)=(y,x)\) acts on \(H_3\) by \(\left(\begin{smallmatrix}0&1\\1&0\end{smallmatrix}\right)\), whose determinant is \(-1\). It therefore disproves that stated image restriction and the claim that the four determinant-one maps generate every component.\par
The needed distinction is the orientation-preserving subgroup. This particular omitted case can be repaired without a new curvature construction: \(\sigma\) is an isometry of the standard equal-sphere product. If \(\phi\) reverses orientation, then \(\psi=\sigma\circ\phi\) preserves orientation and \(\psi^*g_0=\phi^*g_0\). A path for \(\psi\) is already a path with the required endpoint for \(\phi\). Subject to the earlier positive-Ricci construction, this completes the omitted reduction. The defect is Minor; it does not disprove Lemma 1.2.\par
\section*{Coverage, positive checks, and limits}
The group decomposition was checked as a compatible nested sequence of finite cyclic groups. Any finite set of elements of \(\mathbb Q/\mathbb Z\) has a common denominator and generates a cyclic finite subgroup. Compatible generators can be chosen inductively by lifting units in a cyclic quotient. The carry in (17) is needed because the splitting is only a splitting of sets. In (32), exactly one wrap occurs after \(k\) applications of the new generator, giving the old generator on each piece. This identity was independently enumerated for 56 pairs of subgroup order and index, and the general one-wrap proof was checked.\par
Topologically, each punctured \(S^3\times D^4\) piece and every gluing \(S^3\times S^3\) boundary is simply connected. The incidence graph has no cycles, and every finite collection of vertices has a common ancestor because the cyclic subgroups exhaust \(\Gamma\). Van Kampen on finite subtrees, followed by compactness of any loop, supplies simple connectivity of the increasing union. Stabilized metric balls with radii tending to infinity provide completeness of the direct limit; properness of the action additionally uses the separation of the newly inserted orbits, not just freeness. These mechanisms were traced through the inductive conditions and the gluing construction. No purely topological contradiction was established.\par
The base-model Ricci formula (53), the ordinary doubly warped formulas, the circle-bundle correspondence, the \(S^3\)-bundle triviality over \(S^2\), the gauge correction at the endpoint and the order of parameter choices were checked. The explicit map \(\phi_1(u,z)=(u,u^{-1}z)\) does intertwine the two specified actions, and its inverse is \((u,w)\mapsto(u,uw)\). The singular zero section has no bearing on freeness of the weight-one action. The induction for finite or trivial \(\Gamma\) also has direct product/lens-space models and does not require an infinite strictly increasing chain.\par
Several displayed local normalizations need reconciliation in a revision. For example the convex interpolation of \(f\) on p. 40 requires the maximum, rather than the printed minimum, of its two affine functions. The affine tail on p. 25 has cone coordinate \(r+\epsilon10^{-4}/(1-\epsilon)\), so the stated cutoff is not exact. These have explicit local corrections and are not counted again as independent demonstrated obstructions. Small-radius statements must be read with explicit smallness choices; for the constructed product a whole \(S^3_\delta\) cannot lie in a fixed-radius ball when \(\delta\) is arbitrarily large. The differing weight labels in the restatement of Proposition 8.1 likewise need to be reconciled with the actual action in (171). These observations do not substitute for a complete repaired construction.\par
The exact independent execution comprised 24 rational Koszul-curvature cases, four warped-metric specializations, and 56 finite carry-action cases, with all assertions passing. The analytical identities above, rather than a finite parameter grid, establish the findings. Eight source pages were visually checked to disambiguate the key formulas; all source pages were read as text.\par
The full external classification of the mapping class group, all cited compactness results and every smoothing constant in a corrected construction were not independently reproved. In particular this review does not certify a replacement proof after F01–F03. The verdict is version-limited: demonstrated defects in the acquired v3 derivations, with no claim that the 2024 revision or 2025 journal-final article retains them, and no claim that the main theorem is false.\par
\section*{Compact conclusion table}
Item | Judgment | Repair and remaining decision\par
F01, connection/Ricci | Major; asserted mixed term vanishing fails | Reconcile the connection, Yang–Mills condition and metric ansatz. The tested alternative also has a curvature obstruction; no complete repaired construction is certified.\par
F02, radial Ricci | Major; missing second metric derivative | Include the \(g''\) term and redo the neck estimate. The exact warped counterexample refutes the formula, not the final existence theorem.\par
F03, volume normalization | Major; tensor power inconsistent | Use the dimension-six tensor normalization exponent \(-1/3\) and recheck the volume/neck estimates. No full repaired induction is supplied.\par
F04, weight action | Minor; asserted general freeness fails | Restrict to the weight-one action used in the construction; the broader weighted claim has nontrivial isotropy.\par
F05, blow-down description | Minor; incompatible cone statement | Reconcile the universal cone assertion with the displayed \(\mathbb R^3\times S^1\) example. No independent refutation of the construction.\par
F06, mapping-class image | Minor; missing orientation case | Include the determinant-\(-1\) product factor interchange. It supplies the required local correction.\par
Scope and decision | v3 initial review complete | No attribution to the unverified 2024 revision or journal final. Author adoption was not requested; scientific/editorial disposition remains pending.\par
9 E. Bruè, A. Naber and D. Semola, Fundamental groups and the Milnor conjecture, arXiv:2303.15347v3 (20 November 2023), manifest-bound PDF and TeX. https://arxiv.org/abs/2303.15347v3. Annals of Mathematics, article identity and revision metadata, 201 (2025), 225–289. https://annals.math.princeton.edu/2025/201-1/p04.\par
\end{document}
